\documentclass[../../main.tex]{subfiles}

\begin{document}

\section{实数序列}\label{section:ch01-实数序列}

一个实数序列 $\{a_{n}\}_{n=1}^{\infty}$ 是从自然数集到实数的映射. 序列是有界的, 若其象是有界的. 若对所有 $n$, $a_{n}\leqslant a_{n+1}$, 我们说该序列是递增的. 若对所有 $n$, $a_{n}\geqslant a_{n+1}$, 我们说该序列是递减的. 单调指的是或者递增或者递减.

\begin{definition}[收敛]\label{definition:ch01-收敛}
我们说序列 $\{a_{n}\}$ 收敛于实数 $a$, 若对每个 $\varepsilon>0$, 存在指标 $N$ 使得对所有 $n\geqslant N$,
\begin{align*}
|a_{n}-a|<\varepsilon.
\end{align*}
实数 $a$ 称为序列 $\{a_{n}\}$ 的极限.
\end{definition}

\begin{proposition}\label{proposition:ch01-极限性质}
收敛序列的极限是唯一的. 此外, 每个收敛序列是有界的, 且若序列 $\{a_{n}\}$ 与 $\{b_{n}\}$ 分别收敛于 $a$ 与 $b$, 则对所有 $n$, $a_{n}\leqslant b_{n}$ 蕴涵 $a\leqslant b$.
\end{proposition}

\begin{proof}
假定 $\{a_{n}\}$ 既收敛于 $a$ 又收敛于 $a'$. 令 $\varepsilon>0$. 选取指标 $N_{1}$ 和 $N_{2}$ 使得对 $n\geqslant N_{1}$, $|a_{n}-a|<\varepsilon/2$, 而对 $n\geqslant N_{2}$, $|a_{n}-a'|<\varepsilon/2$. 对 $n\geqslant\max\{N_{1},N_{2}\}$, 三角不等式给出
\begin{align*}
|a-a'|\leqslant|a-a_{n}|+|a_{n}-a'|<\varepsilon.
\end{align*}
由于对每个 $\varepsilon>0$, $|a-a'|<\varepsilon$, 我们有 $a=a'$.

为证明序列 $\{a_{n}\}$ 是有界的, 选取 $N$ 使得对 $n\geqslant N$, $|a_{n}-a|<1$. 对 $n\geqslant N$, $|a_{n}|\leqslant|a|+1$. 因此集合 $\{a_{n}\}$ 以 $\max\{|a_{1}|,\cdots,|a_{N}|,|a|+1\}$ 为上界, 以 $-\max\{|a_{1}|,\cdots,|a_{N}|,|a|+1\}$ 为下界.

最后, 假定对所有 $n$, $a_{n}\leqslant b_{n}$, 而 $a>b$. 令 $\varepsilon=(a-b)/2$. 选取 $N$ 使得对 $n\geqslant N$, $|a_{n}-a|<\varepsilon$ 与 $|b_{n}-b|<\varepsilon$. 对这样的 $n$,
\begin{align*}
a_{n}>a-\varepsilon=(a+b)/2,\qquad b_{n}<b+\varepsilon=(a+b)/2,
\end{align*}
因此 $a_{n}>b_{n}$, 与 $a_{n}\leqslant b_{n}$ 矛盾. 因此 $a\leqslant b$.
\end{proof}

\begin{theorem}[单调收敛准则]\label{theorem:ch01-单调收敛}
每个单调有界实数序列收敛于实数.
\end{theorem}

\begin{proof}
令 $\{a_{n}\}$ 为递增有界序列. 则实数集 $E=\{a_{n}\mid n\in\mathbf{N}\}$ 有上界. 根据 $\mathbf{R}$ 的完备性, $E$ 有上确界. 定义 $a=\sup E$. 我们宣称 $\{a_{n}\}$ 收敛于 $a$. 令 $\varepsilon>0$. 由于 $a-\varepsilon$ 不是 $E$ 的上界, 存在指标 $N$ 使得 $a_{N}>a-\varepsilon$. 由于序列是递增的, 对所有 $n\geqslant N$, $a_{n}\geqslant a_{N}>a-\varepsilon$. 另一方面, 对所有 $n$, $a_{n}\leqslant a$, 因此 $|a_{n}-a|<\varepsilon$. 所以 $\{a_{n}\}$ 收敛于 $a$. 递减的情形类似.
\end{proof}

给定序列 $\{a_{n}\}_{n=1}^{\infty}$ 和严格递增的自然数序列 $\{n_{k}\}_{k=1}^{\infty}$, 序列 $\{a_{n_{k}}\}_{k=1}^{\infty}$ 称为 $\{a_{n}\}$ 的子序列.

\begin{theorem}[Bolzano-Weierstrass 定理]\label{theorem:ch01-BW}
每个有界实数序列有收敛的子序列.
\end{theorem}

\begin{proof}
令 $\{a_{n}\}$ 为有界序列. 则其象包含于有界区间 $[m,M]$. 定义闭区间族 $I_{0}=[m,M]$. 若序列的象包含于某个单点集, 则整个序列收敛到那个点. 否则, 将 $I_{0}$ 分为两个等长的闭子区间 $I_{0}^{l}$ 与 $I_{0}^{r}$. 至少有一个这样的子区间包含 $a_{n}$ 对无穷多个 $n$. 选取这样的一个, 称为 $I_{1}$. 归纳地, 对每个自然数 $k$, 得到闭有界区间 $I_{k}$, 长度为 $(M-m)/2^{k}$, 包含 $a_{n}$ 对无穷多个 $n$, 且 $I_{k+1}\subseteq I_{k}$. 根据集套定理（\cref{theorem:ch01-集套}）, 存在点 $x_{0}$ 属于每个 $I_{k}$.

选取指标 $n_{1}$ 使得 $a_{n_{1}}\in I_{1}$. 归纳地, 一旦选取了 $n_{k}$, 选取 $n_{k+1}>n_{k}$ 使得 $a_{n_{k+1}}\in I_{k+1}$. 这是可能的, 因为 $I_{k+1}$ 包含 $a_{n}$ 对无穷多个 $n$, 而只有有限个 $n$ 不超过 $n_{k}$. 则 $x_{0}$ 与 $a_{n_{k+1}}$ 都属于 $I_{k+1}$, 其长度为 $(M-m)/2^{k+1}$. 因此
\begin{align*}
|a_{n_{k+1}}-x_{0}|\leqslant(M-m)/2^{k+1}.
\end{align*}
给定 $\varepsilon>0$, 选取 $K$ 使得 $(M-m)/2^{K+1}<\varepsilon$. 则对 $k\geqslant K$, $|a_{n_{k}}-x_{0}|<\varepsilon$. 因此子序列 $\{a_{n_{k}}\}$ 收敛于 $x_{0}$.
\end{proof}

\begin{definition}[Cauchy 序列]\label{definition:ch01-Cauchy}
实数序列 $\{a_{n}\}$ 称为 Cauchy 序列, 若对每个 $\varepsilon>0$, 存在指标 $N$ 使得对所有 $m,n\geqslant N$,
\begin{align*}
|a_{m}-a_{n}|<\varepsilon.
\end{align*}
\end{definition}

\begin{theorem}[Cauchy 收敛准则]\label{theorem:ch01-Cauchy}
实数序列收敛当且仅当它是 Cauchy 序列.
\end{theorem}

\begin{proof}
首先假定 $\{a_{n}\}$ 收敛于 $a$. 令 $\varepsilon>0$. 选取 $N$ 使得对 $n\geqslant N$, $|a_{n}-a|<\varepsilon/2$. 则对 $m,n\geqslant N$,
\begin{align*}
|a_{m}-a_{n}|\leqslant|a_{m}-a|+|a-a_{n}|<\varepsilon.
\end{align*}
因此序列是 Cauchy 的.

反之, 令 $\{a_{n}\}$ 为 Cauchy 序列. 我们证明它是有界的. 选取 $N$ 使得对 $m,n\geqslant N$, $|a_{m}-a_{n}|<1$. 于是对 $n\geqslant N$, $|a_{n}|\leqslant|a_{N}|+1$. 因此该序列以 $\max\{|a_{1}|,\cdots,|a_{N}|,|a_{N}|+1\}$ 为界. 根据 Bolzano-Weierstrass 定理, 存在收敛子序列 $a_{n_{k}}\to a$. 我们宣称原序列收敛于 $a$. 令 $\varepsilon>0$. 选取 $N_{1}$ 使得对 $m,n\geqslant N_{1}$, $|a_{m}-a_{n}|<\varepsilon/2$. 选取 $K$ 使得对 $k\geqslant K$, $|a_{n_{k}}-a|<\varepsilon/2$, 且 $n_{K}\geqslant N_{1}$. 则对 $n\geqslant N_{1}$, 选取 $k\geqslant K$ 使得 $n_{k}\geqslant n\geqslant N_{1}$, 有
\begin{align*}
|a_{n}-a|\leqslant|a_{n}-a_{n_{k}}|+|a_{n_{k}}-a|<\varepsilon/2+\varepsilon/2=\varepsilon.
\end{align*}
因此 $\{a_{n}\}$ 收敛于 $a$.
\end{proof}

\begin{theorem}\label{theorem:ch01-线性单调}
若序列 $\{a_{n}\}$ 与 $\{b_{n}\}$ 分别收敛于 $a$ 与 $b$, 则对任意实数 $\alpha$ 和 $\beta$, $\{\alpha a_{n}+\beta b_{n}\}$ 收敛于 $\alpha a+\beta b$. 此外, 若对所有 $n$, $a_{n}\leqslant b_{n}$, 则 $a\leqslant b$.
\end{theorem}

\begin{proof}
线性性从三角不等式得到: 对 $\varepsilon>0$, 选取 $N$ 使得对 $n\geqslant N$, $|a_{n}-a|<\varepsilon/(2(|\beta|+1))$ 与 $|b_{n}-b|<\varepsilon/(2(|\alpha|+1))$. 则对 $n\geqslant N$,
\begin{align*}
|\alpha a_{n}+\beta b_{n}-(\alpha a+\beta b)|\leqslant|\alpha||a_{n}-a|+|\beta||b_{n}-b|<\varepsilon.
\end{align*}
单调性已在\cref{proposition:ch01-极限性质}中证明.
\end{proof}

我们说序列 $\{a_{n}\}$ 收敛于 $\infty$, 若对每个实数 $M$, 存在 $N$ 使得对 $n\geqslant N$, $a_{n}\geqslant M$. 类似地, 收敛于 $-\infty$.

\begin{definition}[上极限与下极限]\label{definition:ch01-limsup}
对实数序列 $\{a_{n}\}$, 定义上极限 $\limsup a_{n}$ 与下极限 $\liminf a_{n}$ 为
\begin{align*}
\limsup_{n\to\infty}a_{n}=\inf_{k\geqslant 1}\sup_{n\geqslant k}a_{n},\qquad\liminf_{n\to\infty}a_{n}=\sup_{k\geqslant 1}\inf_{n\geqslant k}a_{n}.
\end{align*}
\end{definition}

\begin{proposition}\label{proposition:ch01-limsup性质}
对实数序列 $\{a_{n}\}$ 和 $\{b_{n}\}$,
\begin{enumerate}[(i)]
\item $\liminf a_{n}\leqslant\limsup a_{n}$, 且序列收敛当且仅当二者相等且为实数.
\item 若对所有 $n$, $a_{n}\leqslant b_{n}$, 则 $\liminf a_{n}\leqslant\liminf b_{n}$ 且 $\limsup a_{n}\leqslant\limsup b_{n}$.
\item 子序列 $\{a_{n_{k}}\}$ 的上极限不超过 $\limsup a_{n}$, 下极限不小于 $\liminf a_{n}$.
\end{enumerate}
\end{proposition}

给定序列 $\{a_{n}\}$, 称序列 $\{s_{n}\}$ 为级数 $\sum_{k=1}^{\infty}a_{k}$ 的部分和序列, 若对每个 $n$, $s_{n}=\sum_{k=1}^{n}a_{k}$. 级数称为可和的, 若其部分和序列收敛.

\begin{proposition}\label{proposition:ch01-级数收敛}
级数 $\sum_{k=1}^{\infty}a_{k}$ 可和当且仅当对每个 $\varepsilon>0$, 存在 $N$ 使得对所有 $m>n\geqslant N$,
\begin{align*}
\left|\sum_{k=n+1}^{m}a_{k}\right|<\varepsilon.
\end{align*}
特别地, 若级数可和, 则 $\lim_{k\to\infty}a_{k}=0$.
\end{proposition}

\begin{proof}
这是 Cauchy 收敛准则应用于部分和序列 $\{s_{n}\}$ 的直接推论.
\end{proof}

\begin{exercise}\label{exercise:ch01-38}
证明: 若 $\{a_{n}\}$ 收敛于 $a$, 则 $\{|a_{n}|\}$ 收敛于 $|a|$. 反过来对吗?
\end{exercise}

\begin{exercise}\label{exercise:ch01-39}
证明: 若 $\{a_{n}\}$ 收敛于 $a$, 则
\begin{align*}
\lim_{n\to\infty}\frac{a_{1}+\cdots+a_{n}}{n}=a.
\end{align*}
反过来对吗?
\end{exercise}

\begin{exercise}\label{exercise:ch01-40}
令 $0<\rho<1$. 证明: 若对所有 $n$, $|a_{n+1}-a_{n}|\leqslant\rho|a_{n}-a_{n-1}|$, 则 $\{a_{n}\}$ 是 Cauchy 序列.
\end{exercise}

\begin{exercise}\label{exercise:ch01-41}
证明 $\lim_{n\to\infty}n^{1/n}=1$.
\end{exercise}

\begin{exercise}\label{exercise:ch01-42}
证明 $\lim_{n\to\infty}(1+1/n)^{n}=e$, 其中 $e=\sum_{k=0}^{\infty}1/k!$.
\end{exercise}

\begin{exercise}\label{exercise:ch01-43}
证明: 若 $\{a_{n}\}$ 是单调的, 则它收敛当且仅当它是有界的.
\end{exercise}

\begin{exercise}\label{exercise:ch01-44}
给出一个有界序列的例子, 它没有收敛的子序列.
\end{exercise}

\begin{exercise}\label{exercise:ch01-45}
证明: 若 $\limsup a_{n}=-\infty$, 则 $\{a_{n}\}$ 收敛于 $-\infty$. 若 $\liminf a_{n}=\infty$, 则 $\{a_{n}\}$ 收敛于 $\infty$.
\end{exercise}

\begin{exercise}\label{exercise:ch01-46}
证明 $\limsup(a_{n}+b_{n})\leqslant\limsup a_{n}+\limsup b_{n}$, 若右边不出现 $\infty-\infty$ 型.
\end{exercise}

\end{document}